\documentclass[11pt,a4paper,oneside]{article}

\usepackage[T1]{fontenc}
\usepackage[utf8]{inputenc}
\usepackage[polish]{babel}
\usepackage{lmodern}
\usepackage{microtype}
\usepackage[a4paper,margin=2.25cm]{geometry}
\usepackage{amsmath,amssymb,mathtools}
\usepackage{booktabs,longtable,tabularx,array}
\usepackage{enumitem}
\usepackage{xcolor}
\usepackage{hyperref}
\usepackage[nameinlink,noabbrev]{cleveref}
\usepackage{tikz}
\usetikzlibrary{arrows.meta,positioning,fit,shapes.geometric}
\usepackage{listings}
\usepackage{fancyhdr}
\usepackage{lastpage}

\definecolor{codegray}{RGB}{245,245,245}
\definecolor{darkblue}{RGB}{20,55,100}
\definecolor{darkgreen}{RGB}{25,105,75}
\definecolor{darkred}{RGB}{145,40,40}

\hypersetup{
  colorlinks=true,
  linkcolor=darkblue,
  urlcolor=darkblue,
  citecolor=darkgreen,
  pdftitle={Dokumentacja techniczna przetwarzania Smog AI 1.7.0},
  pdfauthor={Edward Szczypka}
}

\lstdefinestyle{code}{
  backgroundcolor=\color{codegray},
  basicstyle=\ttfamily\small,
  breaklines=true,
  frame=single,
  columns=fullflexible,
  showstringspaces=false,
  keywordstyle=\color{darkblue},
  commentstyle=\color{darkgreen},
  stringstyle=\color{darkred}
}

\setlength{\headheight}{14pt}
\pagestyle{fancy}
\fancyhf{}
\lhead{Smog AI 1.7.0 --- przetwarzanie}
\rhead{Edward Szczypka}
\cfoot{\thepage/\pageref{LastPage}}

\title{\textbf{Smog AI 1.7.0}\\[0.2em]
\textbf{Dokumentacja techniczna przetwarzania danych, treningu, predykcji i publikacji}}
\author{Edward Szczypka\\\href{mailto:edward@szczypka.guru}{edward@szczypka.guru}}
\date{\today}

\begin{document}
\maketitle

\begin{abstract}
Dokument opisuje kompletny przepływ systemu Smog AI: od pobrania danych GIOŚ
oraz IMGW, przez lokalną bazę SQLite, walidację Pandera, przygotowanie
wersjonowanych zbiorów uczących, otwartą platformę modeli, prognozy godzinowe
na 1--48 godzin, przestrzenną interpolację Polski i publikację do DigitalOcean
Spaces, aż do publicznego FastAPI i dashboardu Streamlit. Dokument jest
techniczny: wskazuje kontrakty danych, kolejność etapów, niezmienniki,
idempotencję, punkty rozszerzeń, obsługę błędów i procedury operacyjne.
Specyfikacja matematyczna modeli znajduje się w odrębnym dokumencie
\texttt{DOKUMENTACJA\_MODELU\_GODZINOWEGO\_PL.tex}.
\end{abstract}

\tableofcontents
\newpage

\section{Granice systemu}

System składa się z dwóch niezależnych komponentów.

\subsection{Komputer lokalny}

Komputer lokalny odpowiada za wszystkie operacje obliczeniowe:
\begin{itemize}[nosep]
  \item pobieranie danych GIOŚ i IMGW;
  \item historię pomiarów w SQLite;
  \item walidację, kontrolę jakości i dopasowanie stacji;
  \item publikację pakietu wejściowego do Spaces;
  \item ponowne pobranie pakietu ze Spaces do treningu;
  \item przygotowanie cech, trening i wybór modeli;
  \item prognozy punktowe dla dokładnych godzin;
  \item interpolację przestrzenną całej Polski;
  \item publikację modeli, metryk, prognoz, map i dokumentacji.
\end{itemize}

\subsection{DigitalOcean}

DigitalOcean Spaces jest prywatnym magazynem obiektowym zgodnym z S3.
DigitalOcean App Platform uruchamia FastAPI i Streamlit. Warstwa zdalna:
\begin{itemize}[nosep]
  \item nie pobiera danych źródłowych;
  \item nie trenuje modeli;
  \item nie uruchamia \texttt{model.predict()};
  \item nie interpoluje przestrzennie;
  \item odczytuje wyłącznie gotowe, lokalnie policzone artefakty.
\end{itemize}

\section{Diagram przepływu}

\begin{figure}[ht]
\centering
\resizebox{\textwidth}{!}{%
\begin{tikzpicture}[
  node distance=9mm and 12mm,
  block/.style={draw, rounded corners, align=center, minimum width=33mm,
                minimum height=10mm, fill=blue!6},
  data/.style={draw, rounded corners, align=center, minimum width=33mm,
               minimum height=10mm, fill=green!8},
  cloud/.style={draw, rounded corners, align=center, minimum width=33mm,
                minimum height=10mm, fill=orange!10},
  arrow/.style={-{Latex[length=2.1mm]}, thick}
]
\node[data] (gios) {GIOŚ JSON-LD};
\node[data, below=of gios] (imgw) {IMGW bieżące + archiwum};
\node[block, right=of gios, yshift=-9mm] (normalize) {Normalizacja\\UTC i kontrakty};
\node[block, right=of normalize] (sqlite) {SQLite WAL\\historia lokalna};
\node[block, right=of sqlite] (quality) {Pandera + jakość\\+ Haversine};
\node[cloud, right=of quality] (bronze) {Spaces\\pakiet Bronze};
\node[block, below=18mm of bronze] (features) {Lokalny odczyt\\i cechy godzinowe};
\node[block, left=of features] (models) {Otwarty rejestr\\modeli};
\node[block, left=of models] (predict) {Prognozy stacyjne\\h=1..48};
\node[block, left=of predict] (spatial) {Interpolacja\\przestrzenna};
\node[cloud, below=18mm of spatial] (serving) {Spaces\\Serving + modele};
\node[block, right=of serving] (api) {FastAPI\\tylko odczyt};
\node[block, right=of api] (ui) {Streamlit\\mapa i dokumentacja};

\draw[arrow] (gios) -- (normalize);
\draw[arrow] (imgw) -- (normalize);
\draw[arrow] (normalize) -- (sqlite);
\draw[arrow] (sqlite) -- (quality);
\draw[arrow] (quality) -- (bronze);
\draw[arrow] (bronze) -- (features);
\draw[arrow] (features) -- (models);
\draw[arrow] (models) -- (predict);
\draw[arrow] (predict) -- (spatial);
\draw[arrow] (spatial) -- (serving);
\draw[arrow] (serving) -- (api);
\draw[arrow] (api) -- (ui);
\end{tikzpicture}}
\caption{Przepływ end-to-end. Strzałka Spaces $\rightarrow$ lokalny trening jest
celowa i realizuje wymaganie komunikacji z Cloud Storage.}
\end{figure}

\section{Konfiguracja i przenośność}

Projekt nie ma zaszytego katalogu instalacji. Katalog repozytorium jest
ustalany na podstawie położenia skryptu albo przez
\texttt{SMOG\_AI\_PROJECT\_ROOT}. Dane uruchomieniowe są przechowywane poza
repozytorium, domyślnie w \texttt{\%ProgramData\%\textbackslash SmogAI}.

Główne pliki:
\begin{lstlisting}[style=code]
C:\ProgramData\SmogAI\config.yaml
C:\ProgramData\SmogAI\smog-ai.env
C:\ProgramData\SmogAI\server-local.env
C:\ProgramData\SmogAI\data\smog.db
C:\ProgramData\SmogAI\logs\...
C:\ProgramData\SmogAI\models\...
\end{lstlisting}

Sekrety nie są przekazywane w argumentach procesu. Plik środowiskowy jest
ładowany jawnie przez skrypty PowerShell i nie trafia do Git.

\section{Pobieranie GIOŚ}

Klient GIOŚ:
\begin{enumerate}
  \item pobiera listę stacji stronicowanym endpointem JSON-LD;
  \item pobiera sensory każdej stacji;
  \item wybiera PM10 i PM2.5;
  \item pobiera bieżącą serię stanowiska;
  \item normalizuje daty do UTC;
  \item zapisuje odpowiedź diagnostyczną w \texttt{raw\_json};
  \item wstawia rekordy idempotentnie przez ograniczenie UNIQUE.
\end{enumerate}

Błąd pojedynczego sensora nie unieważnia całej kolekcji. Niedostępne endpointy
stanowisk historycznych są klasyfikowane jako pominięte lub ostrzeżenia, a
rzeczywiste błędy sieciowe trafiają do \texttt{collection\_errors}.

\section{Pobieranie IMGW}

\subsection{Dane bieżące}

Endpoint SYNOP dostarcza najnowsze obserwacje stacji. Wartość opadu jest
zapisywana z jawnym polem
\texttt{precipitation\_accumulation\_period\_hours}. Domyślna konfiguracja
traktuje ją jako akumulację sześciogodzinną kończącą się w czasie pomiaru.
System nie dzieli takiej wartości przez sześć i nie tworzy sztucznego rozkładu
godzinowego.

\subsection{Oficjalne archiwum terminowe/SYNOP}

Adapter archiwalny:
\begin{enumerate}
  \item wyznacza miesiące objęte skonfigurowanym lookbackiem;
  \item odczytuje listing rocznego katalogu IMGW;
  \item pobiera miesięczne ZIP-y do lokalnego cache;
  \item liczy SHA-256 i sprawdza integralność ZIP;
  \item czyta CSV według dołączonego oficjalnego nagłówka;
  \item normalizuje stację i czas;
  \item zapisuje temperaturę, wilgotność, ciśnienie, wiatr i opad;
  \item zachowuje kody jakości oraz pochodzenie w \texttt{raw\_json};
  \item zapisuje checksum przetworzonego pliku w
        \texttt{application\_state}.
\end{enumerate}

Ponowne uruchomienie pomija plik, którego checksum i status sukcesu nie
uległy zmianie. Kody pól archiwum są konfigurowalne, dzięki czemu adapter można
zastąpić lub dostosować bez zmiany domeny.

\section{Lokalna baza SQLite}

Baza pracuje w trybie WAL z ustawionym \texttt{busy\_timeout}. Migracje są
prowadzone przez Alembic. Najważniejsze niezmienniki:
\begin{itemize}[nosep]
  \item pomiar jest unikalny według źródła, sensora/stacji, parametru i czasu;
  \item prognoza jest unikalna według modelu, stacji, parametru, czasu
        utworzenia, czasu docelowego i horyzontu;
  \item prognoza jest zapisywana przed czasem docelowym;
  \item nie usuwa się podejrzanych pomiarów --- otrzymują flagi jakości.
\end{itemize}

\section{Walidacja Pandera}

Każdy etap ma osobny kontrakt:
\begin{itemize}[nosep]
  \item pomiary GIOŚ;
  \item pomiary IMGW;
  \item pakiet operacyjny;
  \item godzinowe ramki treningowe;
  \item prognozy;
  \item snapshot i powierzchnie przestrzenne.
\end{itemize}

Raport zawiera czas rozpoczęcia i zakończenia, liczbę wierszy, kolumny,
przypadki błędów i checksum. Polityki \texttt{report} oraz \texttt{fail} są
konfigurowalne. Błąd kontraktu treningowego zatrzymuje aktywację modelu.

\section{Warstwy Bronze, Curated i Serving}

\subsection{Bronze}

Pakiet \texttt{operational-data.json.gz} zawiera stacje, sensory, pomiary i
dopasowania. Manifest ma \texttt{run\_id}, zakres danych, liczności, checksum,
status kompletności i listę brakujących elementów.

\subsection{Curated}

Dla każdego celu tworzony jest wersjonowany zbiór godzinowy. Rekord zawiera:
\begin{lstlisting}[style=code]
measurement_time   # origin_time
target_time
horizon_hours
station/source identifiers
lags and rolling features
target-time cyclic features
target
\end{lstlisting}

Warunek czasu jest walidowany:
\[
  \texttt{target\_time}-\texttt{measurement\_time}
  =\texttt{horizon\_hours}.
\]

\subsection{Serving}

Serving zawiera gotowe prognozy stacyjne, powierzchnie mapowe, model cards,
metryki i dokumentację. App Platform korzysta tylko z tej warstwy.

\subsection{Bezpieczeństwo ramek na danych rzeczywistych}

Normalizacja godzinowa zachowuje luki potrzebne do budowy lagów, lecz wiersz
z brakującą wartością bieżącą nie staje się automatycznie obserwacją uczącą.
Dla PM10, PM2.5 i temperatury punkt bazowy musi mieć rzeczywistą wartość
bieżącą. Dla opadu sześciogodzinnego brak bieżącej akumulacji jest dopuszczalny,
ponieważ WO6G jest raportowane tylko w wybranych terminach, a model hurdle
imputuje brakujące cechy bez wymyślania godzinowych wartości opadu.

Modele temperatury i opadu są trenowane na jednej serii dla każdej unikalnej
stacji IMGW. Dopasowanie wielu stacji GIOŚ do tej samej stacji IMGW nie powiela
obserwacji pogodowych. W modelach PM dopasowanie pozostaje częścią budowy cech
dla konkretnej stacji jakości powietrza.

Reprezentacja \texttt{origin $\times$ horyzont} ma konfigurowalny limit
\texttt{maximum\_training\_rows\_per\_target}. Punkty bazowe są wybierane
deterministycznie i proporcjonalnie według stacji przed ekspansją. Wartość
docelowa jest nadal dołączana z pełnej serii po dokładnym
\texttt{target\_time}; limit ogranicza zużycie pamięci, ale nie zmienia
semantyki horyzontu.

\section{Przygotowanie cech}

Cechy są obliczane wyłącznie z danych dostępnych w chwili bazowej. Obejmują:
\begin{itemize}[nosep]
  \item lagi 1, 3, 6, 12 i 24 godziny;
  \item średnie i odchylenia kroczące;
  \item tempo zmian;
  \item składowe kierunku wiatru;
  \item współrzędne i identyfikatory stacji;
  \item dokładny horyzont 1--48 godzin;
  \item sinus/cosinus godziny, dnia tygodnia i roku czasu docelowego.
\end{itemize}

Opad pozostaje wielkością dla jawnego okresu akumulacji. Jeżeli źródło ma
okres 6 h, cecha i cel oznaczają akumulację 6 h kończącą się w
\texttt{target\_time}; jednostka prezentacyjna brzmi \texttt{mm / 6 h}.

\section{Otwarta platforma modeli}

Domena używa Bridge \texttt{ModelProvider}. Provider implementuje:
\begin{lstlisting}[style=code,language=Python]
class ModelProvider(Protocol):
    name: str
    task: ModelTask

    def fit(self, X, y, *, context): ...
    def predict(self, artifact, X, *, context) -> PredictionBundle: ...
    def describe(self, artifact) -> dict: ...
\end{lstlisting}

Provider nie otrzymuje sesji SQLAlchemy, klienta Spaces, FastAPI ani
Streamlit. Może pochodzić z:
\begin{itemize}[nosep]
  \item rejestru metod wbudowanych;
  \item modułu z \texttt{register\_models(registry)};
  \item entry pointu \texttt{smog\_ai.model\_providers};
  \item import stringu \texttt{module:object}.
\end{itemize}

W ten sposób można dołączyć Huber, GAM, splajny, XGBoost, LightGBM, CatBoost,
PyTorch, model grafowy lub własny model bez przebudowy pipeline'u.

\section{Kolejność treningu}

Trening przebiega kaskadowo:
\begin{enumerate}
  \item model temperatury;
  \item model opadu typu hurdle;
  \item chronologiczne prognozy out-of-fold temperatury i opadu;
  \item dołączenie prognoz pogody jako cech PM;
  \item trening PM10 i PM2.5;
  \item porównanie kandydatów z modelami bazowymi;
  \item zapis niezmiennego artefaktu i model card;
  \item atomowa aktualizacja aktywnego wskaźnika.
\end{enumerate}

Cross-fitting chronologiczny zapobiega przekazaniu modelowi PM rzeczywistej
pogody z przyszłości. Przy niedostatecznej historii aktywowany jest jawny
fallback \texttt{insufficient\_history}.

\section{Prognoza godzinowa}

Dla wspólnego \texttt{origin\_time} system tworzy wiersze dla
$h=1,\ldots,48$. Dla każdego wiersza:
\[
  \texttt{target\_time}=\texttt{origin\_time}+h.
\]

Nie istnieje logika wybierania najbliższego spośród 6/12/24 godzin. Dla pełnej
godziny wynik pochodzi z bezpośredniego modelu warunkowanego horyzontem. Dla
czasu pomiędzy pełnymi godzinami można zastosować interpolację liniową lub
PCHIP wyłącznie pomiędzy sąsiednimi prognozami. Ekstrapolacja jest domyślnie
zabroniona.

\section{Interpolacja przestrzenna}

Dla każdej zmiennej i każdego \texttt{target\_time} powstaje osobna
powierzchnia. Interpolator jest kolejnym wymiennym Bridge:
\begin{itemize}[nosep]
  \item IDW --- metoda domyślna;
  \item RBF --- alternatywa;
  \item kolejne implementacje mogą dodać kriging lub model grafowy.
\end{itemize}

Na mapie odróżniane są:
\begin{itemize}[nosep]
  \item dokładne współrzędne miasta;
  \item środek wybranej komórki siatki;
  \item stacje użyte do interpolacji;
  \item odległość do najbliższej stacji;
  \item miara pewności przestrzennej.
\end{itemize}

\section{Publikacja do Spaces}

Obiekty są zapisywane najpierw pod niezmiennym kluczem wersji. Dopiero po
udanym uploadzie i sprawdzeniu checksum aktualizowany jest mały wskaźnik
\texttt{latest.json}. Dzięki temu App Platform nigdy nie zobaczy częściowo
zapisanego pakietu.

Przykładowe prefiksy:
\begin{lstlisting}[style=code]
datasets/bronze/
datasets/curated-hourly/
models-hourly/
metrics/hourly-models/
forecasts/
maps/
documentation/
\end{lstlisting}

\section{FastAPI}

FastAPI uruchamiane na App Platform:
\begin{itemize}[nosep]
  \item interpretuje pytanie tekstowe do struktury;
  \item rozwiązuje nazwę miejsca;
  \item wybiera pakiet o dokładnym \texttt{target\_time};
  \item odczytuje wartość z gotowej powierzchni;
  \item zwraca dokładny punkt, model, jednostkę i niepewność;
  \item udostępnia dokumentację techniczną i matematyczną.
\end{itemize}

Brak pakietu dla czasu docelowego daje jawną odpowiedź o braku prognozy.
Serwer nie wybiera potajemnie najbliższej godziny.

\section{Dashboard}

Dashboard prezentuje:
\begin{itemize}[nosep]
  \item mapę 2D jako tryb domyślny;
  \item opcjonalne, ograniczone logarytmicznie 3D;
  \item nazwy miast jako warstwę nad kolumnami;
  \item PM10, PM2.5, temperaturę, prawdopodobieństwo i sumę opadu;
  \item dokładny punkt zapytania i środek komórki;
  \item profil godzinowy;
  \item modele, metryki i providerów;
  \item dokumentację Markdown i źródła LaTeX.
\end{itemize}

\section{Idempotencja i odporność}

Każdy etap może zostać powtórzony. Idempotencję zapewniają:
\begin{itemize}[nosep]
  \item ograniczenia UNIQUE w SQLite;
  \item checksum archiwów i artefaktów;
  \item niezmienne klucze wersji;
  \item deterministyczne identyfikatory publikacji;
  \item outbox z ponowieniami i exponential backoff;
  \item aplikacyjna dzierżawa oraz mutex Windows;
  \item atomowe wskaźniki aktywnych modeli i powierzchni.
\end{itemize}

\section{Harmonogram Windows}

\begin{longtable}{@{}p{35mm}p{90mm}@{}}
\toprule
Zadanie & Zakres \\
\midrule
Godzinowe & GIOŚ, bieżące IMGW, walidacja, dopasowanie, weryfikacja, prognozy,
powierzchnie, snapshot i publikacja. \\
Dzienne & Uzupełnianie historii, oficjalne archiwum IMGW, raport jakości,
outbox, czyszczenie i backup dzienny. \\
Tygodniowe & Cloud round-trip, ramki uczące, trening, walidacja, aktywacja,
prognozy i publikacja model cards. \\
Miesięczne & Integralność SQLite, backup miesięczny, retencja i raport miejsca. \\
\bottomrule
\end{longtable}

Każde zadanie używa pełnych ścieżek, interpretera z \texttt{.venv}, blokady,
limitu czasu, logu oraz propaguje kod zakończenia.

\section{Backup i odzyskiwanie}

SQLite jest kopiowana przez Online Backup API, a nie przez bezpośrednie
kopiowanie aktywnego pliku WAL. Backup zawiera SHA-256, rozmiar i wynik
\texttt{PRAGMA quick\_check}. Procedura odzyskiwania:
\begin{enumerate}
  \item zatrzymać zadania i procesy API;
  \item zabezpieczyć pliki bazy, WAL i SHM;
  \item sprawdzić backup;
  \item odtworzyć bazę albo uruchomić migracje do pustej bazy;
  \item ponownie uruchomić kolekcję idempotentną;
  \item sprawdzić Spaces i healthcheck.
\end{enumerate}

\section{Monitoring i diagnostyka}

Healthcheck kontroluje bazę, integralność, ostatnią kolekcję, prognozę,
publikację, outbox, modele, wolne miejsce, Spaces i stan zadań. Logi JSON
zawierają czas, poziom, logger, etap i wyjątek. Raporty Pandera i model cards
są publikowane jako niezależne artefakty audytowe.

\section{Kontrakty odbioru}

Wersja jest technicznie ukończona, gdy:
\begin{enumerate}
  \item migracje przechodzą na pustej i istniejącej bazie;
  \item importer GIOŚ i oba adaptery IMGW są idempotentne;
  \item pakiet Bronze jest kompletny i zwalidowany;
  \item trening rzeczywiście odczytuje dane ze Spaces;
  \item rejestr wykrywa metody wbudowane i pluginy;
  \item prognozy istnieją dla godzin 1--48;
  \item każda prognoza ma dokładny czas bazowy i docelowy;
  \item opad ma jawny okres akumulacji;
  \item powierzchnie istnieją dla wszystkich skonfigurowanych celów;
  \item FastAPI nie wykonuje modelu ani interpolacji;
  \item dashboard pokazuje dokładny punkt, modele i dokumentację;
  \item testy, release gate, wheel i ponowne rozpakowanie ZIP przechodzą.
\end{enumerate}

\section{Polecenia operacyjne}

\begin{lstlisting}[style=code]
python -m smog_ai backfill-imgw-archive
python -m smog_ai upload-operational-data
python -m smog_ai build-hourly-features --source object_store
python -m smog_ai list-model-methods
python -m smog_ai train-hourly
python -m smog_ai predict-hourly
python -m smog_ai publish-documentation
python -m smog_ai storage-health
python -m smog_ai hourly-readiness
python -m smog_ai healthcheck --json
\end{lstlisting}

\section{Dokumenty powiązane}

\begin{itemize}[nosep]
  \item \texttt{DOKUMENTACJA\_MODELU\_GODZINOWEGO\_PL.tex} --- model
        matematyczny;
  \item \texttt{TECHNICAL\_PROCESSING\_PL.md} --- wersja platformowa;
  \item \texttt{MODEL\_PLUGIN\_GUIDE\_PL.md} --- tworzenie providerów;
  \item \texttt{CUSTOM\_MODEL\_PROVIDER\_PL.md} --- kompletny przykład pluginu.
\end{itemize}



\section{Ograniczony trening, profile i szybkie douczanie}

Pełne archiwum pozostaje w SQLite oraz opcjonalnie w ObjectStore. Ograniczeniu
podlega wyłącznie ramka używana w pojedynczym treningu. Interfejs
\texttt{TrainingSetPolicy} rozdziela pipeline od konkretnej metody selekcji.
Domyślna implementacja wykonuje okno kroczące, zachowuje obserwacje najnowsze i
rzadkie, a pozostałą część wybiera deterministycznie warstwowo według stacji,
horyzontu, miesiąca i przedziału zmiennej docelowej.

\subsection{Profile quick i full}

\begin{center}
\begin{tabular}{lrr}
\toprule
Właściwość & quick & full\\
\midrule
PM: okno historyczne & 365 dni & 730 dni\\
Pogoda: okno historyczne & 730 dni & 1095 dni\\
Maks. wierszy na cel & 250000 & 600000\\
Maks. wierszy walidacji & 60000 & 120000\\
Horyzonty na origin & 8 & 12\\
Foldy cross-fit & 2 & 4\\
Modele kwantylowe & nie & tak\\
Budżet ścienny & 30 min & 120 min\\
\bottomrule
\end{tabular}
\end{center}

Horyzonty dzielone są na koszyki $[1,6]$, $[7,12]$, $[13,24]$ oraz
$[25,48]$. Dla każdego originu wybierany jest mały deterministyczny podzbiór,
a faza wyboru zmienia się wraz z originem. Wszystkie horyzonty pozostają więc
reprezentowane globalnie bez mnożenia każdej obserwacji przez 48.

\subsection{Wagi}

Waga obserwacji ma postać
\begin{equation}
 w_i = 2^{-a_i/\tau}
       n_{s(i)}^{-1/2}
       n_{h(i)}^{-1/2},
\end{equation}
gdzie $a_i$ oznacza wiek w dniach, $\tau$ okres połowicznego zaniku,
$n_{s(i)}$ liczebność stacji, a $n_{h(i)}$ liczebność horyzontu. Wagi są
normalizowane do średniej jeden i ograniczane do przedziału $[0.1,10]$.

\subsection{Korektor reszt i drift}

Kosztowny champion $f_\theta$ pozostaje niezmieniony pomiędzy pełnymi
treningami. Mały korektor reszt $g_\varphi$ jest aktualizowany metodą
\texttt{partial\_fit} na dojrzałych, zweryfikowanych prognozach:
\begin{align}
 r_i &= y_i-f_\theta(x_i,h_i),\\
 \widehat y_i &= f_\theta(x_i,h_i)+g_\varphi(z_i).
\end{align}
Nowa wersja jest aktywowana tylko wtedy, gdy zmniejsza MAE na późniejszym
fragmencie chronologicznym. Drift jest sygnalizowany przez relatywny wzrost
MAE albo przekroczenie progu biasu.

\subsection{Operacje i postęp}

\begin{lstlisting}[style=code]
python -m smog_ai training-policy-status --profile quick
python -m smog_ai train-hourly --profile quick
python -m smog_ai train-hourly --profile full
python -m smog_ai update-hourly-residuals
python -m smog_ai hourly-drift-status
\end{lstlisting}

Tygodniowy harmonogram używa profilu quick, miesięczny profilu full, a
codzienny wykonuje weryfikację, aktualizację korektora i kontrolę driftu.
Każdy długi etap zapisuje trwały procent, nazwę zadania i ETA.


\section{RangeAwareBackfillBridge}

Przed każdym pobraniem zakres jest ponownie audytowany w SQLite. Niech
$R_p$ oznacza zbiór oczekiwanych slotów parametru $p$, a $D_p$ sloty obecne.
Pobierane są wyłącznie składowe źródła przecinające
\[
 M_p=R_p\setminus D_p.
\]

Dla \texttt{WO6G} zbiór $R_p$ ma kadencję sześciogodzinną, dlatego pięć godzin
pomiędzy kolejnymi sumami nie jest klasyfikowane jako luka. Cache jest
realizowany przez Bridge \texttt{local}, \texttt{object\_store} albo
\texttt{hybrid}. SQLite pozostaje źródłem prawdy dla kompletności.

\begin{lstlisting}[style=code]
python -m smog_ai data-range-audit
python -m smog_ai plan-missing-ranges --audit-package data-ranges.zip
python -m smog_ai fill-missing-ranges --audit-package data-ranges.zip
python -m smog_ai progress --run-type range-backfill --watch
\end{lstlisting}

Po dwóch próbach bez zmniejszenia liczby brakujących slotów akcja otrzymuje
status \texttt{skipped\_no\_progress}; luka pozostaje w raporcie, lecz nie
powoduje nieskończonego ponawiania źródła.


\section{Generyczny rejestr parametrów jakości powietrza}

Centralny \texttt{AirParameterRegistry} rozdziela pięć ról parametru:
\texttt{collect\_current}, \texttt{historical\_backfill},
\texttt{auxiliary\_feature}, \texttt{forecast\_target} i
\texttt{spatial\_surface}. Zachodzi w szczególności
\[
 C_p=1\not\Rightarrow F_p=1,
\]
czyli samo pobieranie parametru nie uruchamia jego modelu.

Definicja zawiera aliasy, jednostkę kanoniczną, kadencję, zakres walidacji,
kod rocznego API, tokeny nazw plików oraz listę providerów. Wbudowany katalog
obejmuje PM10, PM2.5, NO2, SO2, O3, CO, C6H6, NO oraz NOX, lecz po aktualizacji
aktywne domyślnie pozostają tylko dotychczasowe role PM10 i PM2.5.

Dla pomocniczego parametru $q$ tworzy się cechy
\[
 (q_t,q_{t-1},q_{t-6},q_{t-24}),
\]
łączone po stacji i dokładnym czasie origin. Snapshot oraz manifest map
publikują kod, nazwę, aliasy i jednostkę, dzięki czemu parser zapytania może
rozpoznać także parametr dodany konfiguracyjnie.

\begin{lstlisting}[style=code]
python -m smog_ai air-parameter-catalog
python -m smog_ai collect-gios --parameters "NO2,O3"
python -m smog_ai backfill-gios-history --pollutants "NO2,O3"
python -m smog_ai train-hourly --profile quick --targets "NO2"
\end{lstlisting}

\end{document}

\section{HF20: kontrakt 48 godzin serwowania}
Szczegóły zawiera dodatek
\texttt{DODATEK\_TECHNICZNY\_HF20\_TIME\_CONTRACT\_MLOPS\_PL.tex}.
